\paragraph{Extending the original complex multiplication point.}
Once the prescribed generic point lies on a proper arithmetic model,
its integral extension has the exact CM tangent lattice.

\begin{proposition}\label{geo-cm-integral-map-extension}
Let $R$ be a discrete valuation ring with fraction field $F$, and let
$\mathcal Y/R$ be proper. Let $G/\mathcal Y$ be a smooth separated
commutative group scheme of finite type with semiabelian geometric fibres,
equipped with its original rigidified symmetric cubical line $\Theta$ and coefficient
action. Suppose a generic point $x_F:\operatorname{Spec}F\to\mathcal Y$
is supplied, together with an abelian scheme $A/R$ with coefficient action
and a rigidified symmetric cubical line, and a generic
isomorphism $f_F:A_F\simeq x_F^*G$. Assume $f_F$ respects the coefficient
action and a given isomorphism of the original rigidified cubical lines.
Then $x_F$ extends uniquely to an integral point $x_R$, and $f_F$
extends uniquely to an integral isomorphism respecting those same data.
It identifies the exact tangent lattices and their coefficient
eigenkernels.
\end{proposition}
\begin{proof}
The valuative criterion for properness gives $x_R$.
The extension theorem for semiabelian schemes over a normal base
\cite[\S2, Lemma~1]{faltings-semiabelian-extension}
extends both $f_F$ and its inverse. Uniqueness makes their composites
the identity, and also proves compatibility with every coefficient
endomorphism.

Transport the pulled-back line to $A$. The given generic line
isomorphism is a rational section of its ratio with the prescribed line
on $A$. Since $A$ is regular and its special fibre is geometrically
integral and reduced, its divisor is $mA_s$ for some integer $m$.
Pullback along zero gives $m[s]$. The rigidification makes this
rational section equal to $1$ at zero, so $m=0$. Thus the original
generic isomorphism extends to a nowhere vanishing section, with no
vertical twist or change of scalar. Use adjunctions normalized at the
origin; properness of $A_F$ makes that normalization unique.
The cubical and symmetry comparisons, and these normalized adjunction
comparisons, agree generically and hence everywhere by flatness.
Differentiating the resulting group isomorphism identifies the integral
Lie modules and all kernels of their coefficient eigenvalue equations.
\end{proof}

Apply the proposition after the finite extension used to obtain the
good CM model in Proposition~\ref{geo-cm-integral-product}.
The integral point stays in the closure of its selected generic
component, since the ideal of that closure vanishes generically and
$R\hookrightarrow F$. Its tangent eigenlattices are therefore the
same lattices computed by the product isogeny and the different.
Together with Proposition~\ref{geo-ramified-hodge-determinant}, this
supplies the integral comparison used in
Corollary~\ref{geo-flat-norm-assembly} whenever the prescribed generic
CM point and its original-line identification are available.
Proposition~\ref{geo-cm-period-point-descent} supplies that exact generic
identification from moduli point separation and the prescribed pure
periods. The arithmetic model and its pointwise period presentation
remain required.
